\SourceLeaf{372}%
\begin{center}
{\large NACHLASS UND BRIEFWECHSEL.\par}
\bigskip
PRAECEPTA GENERALISSIMA\\
PRO INVENIENDIS CENTRIS CIRCULI OSCULANTIS AD QUODVIS CURVAE DATAE PUNCTUM DATUM.
\end{center}

Definiatur situs puncti cuiusvis in plano dato per duas variabiles $x,y$ (sive sint distantiae a punctis sive lineis datis sive anguli etc.), sintque
\[
P=0,\qquad P'=0
\]
aequationes duarum curvarum, designantibus $P,P'$ functiones quascunque ipsarum $x,y$; porro sit
\[
\begin{aligned}
dP&=p\,dx+q\,dy,&dP'&=p'\,dx+q'\,dy,\\
d^2P&=r\,dx^2+2s\,dx\,dy+t\,dy^2,&d^2P'&=r'\,dx^2+2s'\,dx\,dy+t'\,dy^2.
\end{aligned}
\]
Tunc curvae punctum per
\[
x=\xi,\qquad y=\eta
\]
definitum commune habebunt, si pro hisce valoribus
\[
P=P';
\]
praeterea in hoc puncto eandem tangentem habebunt, si insuper
\[
pq'=p'q;
\]
